\section{Theory: what \rrebel{}-std optimizes, and why}
\label{sec:theory}

This section develops the theory behind \rrebel{} with group standardization
(``-std''), the variant family that wins every comparison in
Section~\ref{sec:results}. The approach follows DPO~\citep{rafailov2023dpo}:
start from the KL-regularized objective, use its closed-form optimum to turn a
statement about rewards into a statement about log-ratios, and then ask what a
given loss on those log-ratios converges to. We answer five questions.
\emph{What is the target?} (\S\ref{sec:th-target}): for \emph{any} discrepancy
$d$ and \emph{any} full-support sampler, the population minimizer is a
KL-regularized optimum; standardization changes only its temperature.
\emph{What does standardization buy?} (\S\ref{sec:th-std}): a trust region
whose radius does not depend on the reward's scale, which matters because in
this problem that scale varies by orders of magnitude across $\lmb$.
\emph{Why is a robust $d$ free?} (\S\ref{sec:th-robust}): differencing within
a pair symmetrizes reward noise, so robust losses do not bias the target.
\emph{How does it relate to GRPO?} (\S\ref{sec:th-grpo}): GRPO without a KL
term is \emph{exactly} \rrebel{}-$\ell_2$-std with the reference refreshed every
step. \emph{Which constraint should the reward impose?}
(\S\ref{sec:th-constraint}): the tradeoff curve plots expectations, and a floor
on every sample does not target them. A floor on the policy's expected
content does, and its optimum puts each floor's point on the line
$\mathbb{E}[c] = \tau$.

Proofs are in Appendix~\ref{app:proofs}. Every statement is also checked
numerically, by a script that runs the repository's actual loss code on a
tabular policy (\texttt{analysis/\allowbreak theory\_checks.py};
Appendix~\ref{app:checks}).

\paragraph{Setting.}
A \emph{context} is $c = (x, \lmb)$, drawn as in training
($x$ from the data, $\lmb \sim \mathcal{U}[0,1]$). Prompts range over a finite
set $\mathcal{Z}$ (all length-$L$ token strings). We analyse the
population reward $r(c,z) = \mathbb{E}_{y}[\Rlam(z,x;y)]$, of which the Monte
Carlo average $\hat r_\lmb$ in Eq.~\eqref{eq:mcreward} is an unbiased estimate.
$\pi_0(\cdot \mid c)$ denotes the reference policy, assumed to have full support
on $\mathcal{Z}$. Everything below holds context by context, so the policy is
written $\pi(\cdot\mid c)$ throughout; the $\lmb$-conditioned network of
Section~\ref{sec:setup} is one way of representing the whole family
$\{\pi(\cdot\mid c)\}_c$ with shared parameters.

% ---------------------------------------------------------------------------
\subsection{The target: reward-difference regression recovers the KL-regularized optimum}
\label{sec:th-target}

Fix a temperature $\tau_c > 0$ for each context and consider
\begin{equation}
  \max_{\pi}\;
  \mathbb{E}_{z\sim\pi(\cdot\mid c)}\bigl[r(c,z)\bigr]
  \;-\; \tau_c\, \mathrm{KL}\bigl(\pi(\cdot\mid c)\,\|\,\pi_0(\cdot\mid c)\bigr).
  \label{eq:th-klobj}
\end{equation}

\begin{lemma}[KL-regularized optimum; \citealp{peters2007rwr,rafailov2023dpo}]
\label{lem:gibbs}
The unique maximizer of Eq.~\eqref{eq:th-klobj} is the Gibbs policy
$\pi^\star_{\tau_c}(z\mid c) = \pi_0(z\mid c)\,e^{r(c,z)/\tau_c}/Z(c)$, and for
all $z, z'$
\begin{equation}
  r(c,z) - r(c,z') \;=\;
  \tau_c\Bigl[\log\tfrac{\pi^\star_{\tau_c}(z\mid c)}{\pi_0(z\mid c)}
            - \log\tfrac{\pi^\star_{\tau_c}(z'\mid c)}{\pi_0(z'\mid c)}\Bigr].
  \label{eq:th-diff}
\end{equation}
\end{lemma}

The partition function cancels in differences, which is the same cancellation
DPO exploits. DPO uses it to express a Bradley--Terry preference
\emph{probability} through log-ratio differences and fits it by
classification. REBEL~\citep{gao2024rebel} and \rrebel{} use it to express a
reward \emph{difference} through log-ratio differences, and fit it by
regression.

\paragraph{The population loss.}
Let $q(\cdot\mid c)$ be the distribution prompts are sampled from, let
$s(c) > 0$ be a reward scale, and let $d$ be any discrepancy with $d(0)=0$ and
$d(e) > 0$ for $e \neq 0$. The population version of Eq.~\eqref{eq:rrebel} is
\begin{equation}
  \mathcal{L}_{q,s}(\pi) =
  \mathbb{E}_c\,\mathbb{E}_{z,z'\sim q(\cdot\mid c)}\,
  d\Bigl(\beta\bigl[\log\tfrac{\pi(z\mid c)}{\pi_0(z\mid c)}
                   - \log\tfrac{\pi(z'\mid c)}{\pi_0(z'\mid c)}\bigr]
         \;-\; \tfrac{r(c,z) - r(c,z')}{s(c)}\Bigr).
  \label{eq:th-poploss}
\end{equation}
Plain \rrebel{} has $s \equiv 1$. \rrebel{}-std has
$s(c) = \sigma_q(c)$, the standard deviation of $r(c,\cdot)$ under
$q(\cdot\mid c)$: the population counterpart of the within-group standard
deviation in Eq.~\eqref{eq:rewardprep}.

\begin{theorem}[Consistency for every discrepancy and every full-support sampler]
\label{thm:consistency}
Suppose $q(\cdot\mid c)$ has full support on $\mathcal{Z}$ for almost every $c$.
Then $\mathcal{L}_{q,s}(\pi) \ge 0$, with equality if and only if
$\pi(\cdot\mid c) = \pi^\star_{\beta s(c)}(\cdot\mid c)$ for almost every $c$.
In particular:
\begin{enumerate}[label=(\roman*),itemsep=1pt,topsep=2pt]
  \item plain \rrebel{} is minimized uniquely by the KL-regularized optimum at
        temperature $\beta$, whatever the discrepancy $d$;
  \item \rrebel{}-std is minimized uniquely by the KL-regularized optimum at the
        \emph{context-dependent} temperature $\beta\,\sigma_q(c)$. Equivalently,
        it solves KL-regularized RL on per-context \emph{standardized} reward,
        \[
          \max_\pi\; \mathbb{E}_c\bigl[\mathbb{E}_{\pi}[\,r/\sigma_q(c)\,]
          - \beta\,\mathrm{KL}(\pi\,\|\,\pi_0)\bigr].
        \]
\end{enumerate}
\end{theorem}

Three consequences are worth stating.
\emph{First, the target does not depend on how prompts are sampled.} The
perturbed samplers used for exploration in the original \rrebel{}
formulation (e.g.\ excluding the top-$m$ tokens) change the loss surface but
not its minimizer, so exploration costs no bias. For \rrebel{}-std the
qualification is that $\sigma_q$ is itself measured under $q$; we analyse the
natural choice below, $q = \pi_0$ at the start of each reference window.
\emph{Second, the discrepancy enters only through ``zero iff the residual is
zero''.} At the population level, $\ell_1$, Huber, and Cauchy losses all have the
same target as $\ell_2$; \S\ref{sec:th-robust} shows that this survives noisy
rewards.
\emph{Third, amortization is lossless.} Because the loss vanishes if and only if
it vanishes context by context, a $\lmb$-conditioned policy class rich enough to
represent every $\pi^\star(\cdot\mid x,\lmb)$ recovers all of them
simultaneously. A single policy can trace the whole frontier without giving up
any point on it. (Like DPO's Theorem~1, this is a statement about the
nonparametric minimizer; with finite capacity, which contexts the fit favours
is the subject of Proposition~\ref{prop:balance}.)

\begin{lemma}[Standardization enlarges the reward equivalence class]
\label{lem:invariance}
DPO calls two rewards equivalent if they differ by a function of the context
alone, $r' = r + b(c)$; such rewards share an optimal policy. Plain \rrebel{}
has exactly these classes. Under \rrebel{}-std, every
$r' = a(c)\,r + b(c)$ with $a(c) > 0$ yields the \emph{same loss}
$\mathcal{L}_{q,\sigma_q}$ for every $\pi$, not merely the same minimizer.
\end{lemma}

The constants $\kappa$ and $g(\lmb)$ in Eq.~\eqref{eq:rewardprep} are such an
$a(c)$. Lemma~\ref{lem:invariance} is therefore the formal version of the
earlier observation that they have no effect on any -std variant. Hand-tuning
a $\lmb$-dependent reward weight is unnecessary: standardization already puts
every context on a common scale.

% ---------------------------------------------------------------------------
\subsection{What standardization buys: a scale-free trust region}
\label{sec:th-std}

Theorem~\ref{thm:consistency}(ii) says standardization replaces the global
temperature $\beta$ with $\beta\sigma(c)$. The next three results explain why
this is the right temperature for a $\lmb$-conditioned policy. Throughout, the
sampler at the start of a reference window is the reference itself,
$q = \pi_0$. This holds exactly at every refresh and approximately during the
$K = 150$ steps that follow.

\begin{proposition}[Scale-free trust region]
\label{prop:trust}
Fix $c$ and let $u = (r - \mathbb{E}_{\pi_0}r)/\sigma_{\pi_0}(c)$ be the
standardized reward, with cumulant generating function
$\Lambda_u(t) = \log\mathbb{E}_{\pi_0}e^{tu}$. The \rrebel{}-std target
$\pi^+ = \pi^\star_{\beta\sigma}$ satisfies
\begin{equation}
  \mathrm{KL}(\pi^+\|\pi_0) = \tfrac{1}{\beta}\Lambda_u'\!\bigl(\tfrac1\beta\bigr)
                              - \Lambda_u\!\bigl(\tfrac1\beta\bigr),
  \qquad
  \mathbb{E}_{\pi^+}[u] = \Lambda_u'\!\bigl(\tfrac1\beta\bigr).
  \label{eq:th-trust}
\end{equation}
Both quantities depend on the context only through the \emph{shape} of the law
of $u$, and not through the reward's scale or location. Moreover $\pi^+$ solves
$\max_\pi \mathbb{E}_\pi[r]$ subject to
$\mathrm{KL}(\pi\|\pi_0) \le \varepsilon(c)$, where $\varepsilon(c)$ is the
left-hand side of Eq.~\eqref{eq:th-trust}. For large $\beta$,
\begin{equation}
  \varepsilon(c) = \tfrac{1}{2\beta^2} + \tfrac{\gamma(c)}{3\beta^3} + O(\beta^{-4}),
  \qquad
  \mathbb{E}_{\pi^+}[u] = \tfrac{1}{\beta} + \tfrac{\gamma(c)}{2\beta^2} + O(\beta^{-3}),
\end{equation}
where $\gamma(c) = \mathbb{E}_{\pi_0}[u^3]$ is the reward's skewness. The plain
target $\pi^\star_\beta$ has instead
$\mathrm{KL}(\pi^\star_\beta\|\pi_0) = \sigma(c)^2/(2\beta^2) + O(\sigma^3/\beta^3)$.
\end{proposition}

Read as an algorithm, each reference window of \rrebel{}-std is a
\emph{trust-region} step. To leading order the KL radius is the same in every
context and the expected gain is $1/\beta$ standard deviations. For
Gaussian-shaped rewards the radius is exactly $1/(2\beta^2)$, i.e.\ $2$ nats at
our $\beta = 0.5$. Plain \rrebel{} applies one KL \emph{penalty}, so its radius
grows as $\sigma(c)^2$. Figure~\ref{fig:theory-toy} (left) shows the
consequence on a toy problem: over five decades of reward scale, the
standardized target's KL is constant, while the plain target's KL ranges from
$10^{-6}$ nats (no movement) to $\log|\mathcal{Z}|$ (collapse to a point mass
in a single window).

\begin{proposition}[A lagged reference is mirror ascent with standardized steps]
\label{prop:mirror}
Suppose each reference window is solved exactly, and the reference is then
refreshed to the result. The iterates satisfy
\begin{equation}
  \pi_{k+1}(z\mid c) \;\propto\; \pi_k(z\mid c)\,
  \exp\bigl(\eta_k(c)\, r(c,z)\bigr),
  \qquad \eta_k(c) = \frac{1}{\beta\,\sigma_{\pi_k}(c)},
  \label{eq:th-mirror}
\end{equation}
which is mirror ascent on $\mathbb{E}_\pi[r]$ with the KL mirror map
\citep{beck2003mirror}, i.e.\ exponentiated-gradient / natural policy gradient
\citep{kakade2001npg}. This is the connection REBEL draws for its own update
\citep{gao2024rebel}, here with a per-context step size. Unrolled,
$\pi_k \propto \pi_0\, e^{T_k(c)\, r}$ with $T_k = \sum_{j<k}\eta_j$. Hence every
iterate is a Gibbs distribution: no single window can collapse the policy. As
$T_k \to \infty$ the iterates concentrate on $\arg\max_z r(c,z)$, the best
prompt for that $(x,\lmb)$, which is the operating point the multi-objective
problem asks for.
\end{proposition}

With standardization, the step $\eta_k$ \emph{grows} as the policy sharpens and
$\sigma_{\pi_k}$ shrinks, which holds the per-window KL near
$\varepsilon \approx 1/(2\beta^2)$. Progress therefore does not stall as the
policy concentrates, and it is \emph{identical} across contexts whose rewards
differ only in scale. Without standardization, $T_k = k/\beta$ in every
context, so contexts with small reward spread barely move while contexts with
large spread saturate in one window. Figure~\ref{fig:theory-toy} (right) shows
both behaviours. The earlier informal claim that ``\rrebel{}'s optimum is a
distribution'' should be read in this precise sense: each window targets a
Gibbs policy at bounded KL, so concentration happens at a controlled rate of
roughly $1/(2\beta^2)$ nats per window rather than in one step.

\begin{proposition}[Standardization balances the amortized objective]
\label{prop:balance}
At the start of a window ($\pi = \pi_0 = q$, all log-ratios zero), the expected
$\ell_2$ loss of a context is $2\,\sigma(c)^2$ for plain \rrebel{}, and exactly
$2$ for \rrebel{}-std. Under $\ell_1$ it is $\sigma(c)\cdot m_u$ versus $m_u$,
where $m_u = \mathbb{E}|u-u'|$ depends only on the shape of $u$.
\end{proposition}

A shared-parameter policy trained without standardization therefore gives
contexts weight in proportion to $\sigma(c)^2$ (or $\sigma(c)$). It spends its
capacity on the high-variance end of the $\lmb$ range and treats contexts
where every sampled prompt has nearly the same reward as noise. With
standardization, every context starts from the same loss.

\paragraph{How large is the scale variation in this problem?}
Reward~I (Eq.~\ref{eq:reward}) has two branches on very
different scales. Where the content floor is met, the reward is a sentiment
score on $[0,100]$. Where it is not, the reward is $0.01\times$ the shortfall,
typically $|r| \lesssim 0.5$. Which branch
dominates a group depends on both $\lmb$ and $x$.
Section~\ref{sec:res-scale} measures $\sigma(x,\lmb)$ under the training
sampler: it varies by a factor of \sRangeQ{} across contexts. It also computes
the per-group KL radii that Proposition~\ref{prop:trust} predicts for the plain
and standardized targets.

% ---------------------------------------------------------------------------
\subsection{Why a robust discrepancy costs nothing}
\label{sec:th-robust}

\begin{proposition}[Pairwise differencing symmetrizes reward noise]
\label{prop:symm}
Suppose observed rewards are $\hat r(z_i) = r(z_i) + \epsilon_i$, with the
$\epsilon_i$ independent across the members of a group and \emph{either}
(a) identically distributed, of any shape (skewed or heavy-tailed), \emph{or}
(b) each symmetric about zero, with any prompt-dependent scale. Then
$\epsilon_i - \epsilon_j$ is symmetric about zero, and for every even, convex
$d$ the minimizer over $h$ of
$\mathbb{E}\, d\bigl(h - (\hat r(z_i) - \hat r(z_j))/s\bigr)$ is the noiseless
target $(r(z_i) - r(z_j))/s$. It is unique if $d$ is strictly convex, or, for
$d = |\cdot|$, if $\epsilon_i - \epsilon_j$ has positive density at zero.
\end{proposition}

By contrast, a \emph{pointwise} robust regression of a score $h(z)$ onto
$\hat r(z)$ estimates $\mathrm{median}(r + \epsilon)$, which is biased whenever
the noise is skewed. The pairwise construction is what makes the ``R'' in
\rrebel{} free: $\ell_1$ and Huber inherit the target of
Theorem~\ref{thm:consistency}, and their bounded influence functions
\citep{huber1964robust} cap how much any single noisy pair can move the
gradient. This matters here because $\hat r_\lmb$ is an average of a
\emph{discontinuous} per-sample reward: a rewrite whose content lands just
either side of $100\lmb$ switches between a sentiment score and a small
penalty. (Non-convex $d$, such as Cauchy, additionally needs the noise
difference to be unimodal.)

\begin{lemma}[Standardized targets are bounded]
\label{lem:bounded}
With the sample standard deviation of Eq.~\eqref{eq:rewardprep}
($G-1$ normalization), every standardized target satisfies
$|\tilde r_i - \tilde r_j| \le \sqrt{2(G-1)}$, and this bound is attained
\citep{samuelson1968deviant}.
\end{lemma}

For $G = 16$, every regression target lies in $[-5.48,\, 5.48]$, whatever the
reward scale. With Huber's $\delta = 1$, the targets therefore span both its
quadratic and its linear regimes. Together with a bounded influence function,
Lemma~\ref{lem:bounded} means no single pair can dominate a group's update.

% ---------------------------------------------------------------------------
\subsection{Relation to GRPO: GRPO is \rrebel{}-$\ell_2$-std with lag one}
\label{sec:th-grpo}

\begin{proposition}[First-order equivalence]
\label{prop:firstorder}
Let $\psi = d'$ (odd, since $d$ is even). At the start of a reference window,
where $\pol = \polref$ and every log-ratio is zero, the \rrebel{} gradient over
$S$ groups of size $G$ with $P = G(G-1)/2$ pairs is
\begin{equation}
  \nabla_\theta \mathcal{L}^{\rrmath}
  = -\frac{\beta}{S P}\sum_{\text{groups}}\sum_{i=1}^{G}
     u_i\,\nabla_\theta \log\pol(z_i),
  \qquad
  u_i = \sum_{j\neq i}\psi\bigl(\tilde r_i - \tilde r_j\bigr).
  \label{eq:th-firstorder}
\end{equation}
This is a policy gradient with a zero-sum advantage built from pairwise
comparisons, and its shape is the influence function of $d$:
\begin{itemize}[itemsep=1pt,topsep=2pt]
  \item $\ell_2$: $u_i = 2G(\tilde r_i - \overline{\tilde r}) = 2G A_i$, with
        $A_i$ the GRPO advantage of Eq.~\eqref{eq:grpo-adv};
  \item $\ell_1$: $u_i = \sum_{j}\mathrm{sign}(r_i - r_j) = 2\,\mathrm{rank}_i - (G+1)$,
        a \emph{centered-rank} advantage. It is invariant to every strictly
        increasing transformation of the reward, and it is the fitness shaping
        of natural evolution strategies \citep{wierstra2014nes,salimans2017es};
  \item Huber: $u_i = \sum_j \mathrm{clip}(\tilde r_i - \tilde r_j, -\delta, \delta)$,
        which interpolates between the rank ($\delta\to0$) and GRPO
        ($\delta\to\infty$) advantages.
\end{itemize}
\end{proposition}

\begin{corollary}[GRPO as a special case]
\label{cor:grpo}
If the reference is refreshed at every step ($K=1$), every gradient is taken at
a window start. \rrebel{}-$\ell_2$-std then produces, at every step, a positive
constant ($4\beta G L/(G-1)$ for prompts of length $L$) times the gradient of
single-update GRPO with $\beta_{\mathrm{KL}} = 0$. Adam is invariant to a
constant rescaling of the gradient (up to its $\epsilon$, and provided
gradient-norm clipping is inactive), so under Adam the two produce the same
parameter trajectory.
\end{corollary}

Within a window ($K > 1$), the advantage is evaluated at the regression
\emph{residual},
$u_i = \sum_j \psi\bigl(\tilde r_i - \tilde r_j - [\logratio(z_i) - \logratio(z_j)]\bigr)$.
It therefore shrinks as the log-ratio differences absorb the reward differences,
and it vanishes at the Gibbs target of Theorem~\ref{thm:consistency}.
GRPO's advantage $A_i$ does not depend on $\theta$; only its KL term opposes it,
and in this single-update pipeline that term is the \emph{only} regularizer
(Section~\ref{sec:grpo}). Corollary~\ref{cor:grpo} thus locates the difference
between the two families precisely. Relative to GRPO, \rrebel{}-std adds
(1) a \emph{lag} $K$, which gives every window a finite, trust-region target
(Propositions~\ref{prop:trust}--\ref{prop:mirror}), and (2) a
\emph{discrepancy} $d$, whose influence function becomes the advantage shape
(Proposition~\ref{prop:firstorder}), with no change of target
(Proposition~\ref{prop:symm}). The winning variants use $K = 150$ and a
bounded-influence $d$.

\paragraph{Does the choice between $\ell_1$ and Huber matter for settling?}
The decay argument above needs $\psi(e) \to 0$ as $e \to 0$. This holds for
$\ell_2$, Huber and Cauchy, but not for $\ell_1$, whose force on an unfitted pair
is $\pm1$ however small the residual. In the deterministic population limit,
fixed-step Adam on the $\ell_1$ loss therefore stalls at a small error floor,
while Huber converges; in our toy the floor is $8\times10^{-9}$ nats against
$3\times10^{-11}$. With sampled groups, in-group standardization and noisy
rewards, however, the distinction disappears: the stationary jitter of the two
fits is indistinguishable ($5.1\times10^{-5}$ versus $6.5\times10^{-5}$ nats), because
gradient noise dominates. This property therefore does \emph{not} explain the
late-training divergence between the two -std variants reported in
Section~\ref{sec:res-dynamics}, and we do not claim a mechanism for it.

\begin{figure}[t]
\centering
\includegraphics[width=\textwidth]{theory_toy.pdf}
\caption{Propositions~\ref{prop:trust} and~\ref{prop:mirror} on a toy problem
($|\mathcal{Z}| = 400$ prompts in the left panel, $200$ in the right, Gaussian-shaped rewards,
$\beta = 0.5$), computed exactly. \textbf{Left:} KL from the reference to the
one-window target, for rewards of identical shape but different scale
$\sigma$. The standardized target moves the same distance at every scale
($1.82$ nats here). The plain target ranges from negligible movement to
collapse onto a point mass. \textbf{Right:} mass on the best prompt across
refresh windows. The three standardized trajectories ($\sigma = 0.01, 1, 100$)
coincide exactly. Plain \rrebel{} saturates in one window at $\sigma = 100$ and
barely moves in $40$ windows at $\sigma = 0.01$.}
\label{fig:theory-toy}
\end{figure}

% ---------------------------------------------------------------------------
\subsection{Which constraint should the reward impose?}
\label{sec:th-constraint}

The results so far concern how \rrebel{}-std optimizes a given reward. This
subsection asks which reward the tradeoff curve of Eq.~\eqref{eq:curve}
calls for. Fix a context $(x,\lmb)$ and its floor $\tau = 100\lmb$, and
suppress them from the notation. Let $\mathcal{M} \subseteq \mathcal{Z}$ be a
finite set of prompts with expected content $\bar c_z$ and sentiment $\bar s_z$,
and for a distribution $\pi$ on $\mathcal{M}$ write
$C(\pi) = \mathbb{E}_\pi[\bar c_z]$ and $S(\pi) = \mathbb{E}_\pi[\bar s_z]$.
The curve's point for floor $\tau$ is the average over sentences of
$(C(\pi), S(\pi))$.

\begin{lemma}[A deterministic choice from a few prompts clusters the curve]
\label{lem:cluster}
Suppose the policy is deterministic, emitting $z^\star(x,\lmb) \in \mathcal{M}$
with $|\mathcal{M}| = K$. If $z^\star$ does not depend on $x$, the curve has at
most $K$ distinct points. If $z^\star(x,\lmb) = \arg\max_{z\in\mathcal{M}}
r_\lmb(x,z)$ for a reward continuous in $\lmb$, the assignment is piecewise
constant in $\lmb$: two floors share a point unless some
$r_\lmb(x,z) - r_\lmb(x,z')$ changes sign between them.
\end{lemma}

For Reward~I, $r_\lmb(z) = \mathbb{E}_z[s\,\mathbf{1}\{c\ge\tau\}] -
0.01\,\mathbb{E}_z[(\tau - c)_+]$. A crossing happens where the
feasibility-weighted sentiments of two prompts cross, which depends on how each
prompt's \emph{per-sample} content is spread around $\tau$, not on where its
mean lies. Nothing places the crossings at $\bar c = \tau$. The trained
policies of Section~\ref{sec:res-diagnosis} are in exactly this regime.

\begin{proposition}[The $\varepsilon$-constraint over distributions]
\label{prop:eps}
Let $c_0$ be the largest content among the prompts of maximal sentiment and
$c_{\max} = \max_z \bar c_z$. For $\tau \le c_{\max}$ the value
$V(\tau) = \max\{S(\pi) : C(\pi) \ge \tau\}$ equals the upper concave envelope
of the points $\{(\bar c_z, \bar s_z)\}$ evaluated at $\max(\tau, c_0)$. An
optimal $\pi$ mixes at most two prompts. For every binding floor,
$\tau \in (c_0, c_{\max}]$, \emph{every} optimal $\pi$ has $C(\pi) = \tau$
exactly, and $V$ is concave, continuous and non-increasing there.
\end{proposition}

So the curve that Eq.~\eqref{eq:epsconstraint} calls for has the properties one
would want of a tradeoff curve: each binding floor has its own point, and that
point lies on the line $\mathbb{E}[c] = \tau$. Floors below $c_0$ are inactive and
correctly share the point of the most positive policy. The price is that the
optimum generally \emph{randomizes} between two prompts whenever $\tau$ falls
between their contents. A deterministic policy attains $V$ only at the
envelope's vertices, and Lemma~\ref{lem:cluster} applies to it.

\begin{proposition}[What Reward~II converges to]
\label{prop:penalty}
Let $F(\pi) = S(\pi) - \frac{\nu}{2}\bigl(\tau - C(\pi)\bigr)_+^2$, the
per-context objective of Eq.~\eqref{eq:penalty}, and let $\pi_\mu$ denote the
Gibbs policy $\pi_0\,e^{(\bar s + \mu \bar c)/\beta}/Z$.
\begin{enumerate}[label=(\roman*),itemsep=1pt,topsep=2pt]
  \item \emph{One window.} $F(\pi) - \beta\,\mathrm{KL}(\pi\,\|\,\pi_0)$ has a
        unique maximizer, $\pi_{\mu_\nu}$, where $\mu_\nu$ is the unique solution
        of $\mu = \nu\,(\tau - C(\pi_\mu))_+$. Its shortfall is
        $(\tau - C(\pi_{\mu_\nu}))_+ = \mu_\nu/\nu \le \mu^\star/\nu$, where
        $\mu^\star$ is the multiplier of the KL-regularized constrained problem
        ($C(\pi_{\mu^\star}) = \tau$). The shortfall therefore vanishes as
        $\nu \to \infty$.
  \item \emph{Across windows.} The maximizers of $F$ itself mix at most two
        prompts and have $C(\pi) = t^\star$ with
        $\tau - t^\star = |U'(t^\star)|/\nu$ for $\tau > c_0$, where $U$ is the
        upper concave envelope of Proposition~\ref{prop:eps} and $U'$ its slope.
  \item \emph{The reward.} The functional gradient of $F$ at $\pi$ is
        $\bar s_z + \mu\,\bar c_z$ with $\mu = \nu(\tau - C(\pi))_+$. Reward~II
        (Eq.~\ref{eq:reward-exp}) is this gradient with $C(\pi)$ estimated by the
        group mean, up to a per-group affine map that standardization removes
        (Lemma~\ref{lem:invariance}).
\end{enumerate}
\end{proposition}

By Theorem~\ref{thm:consistency}, each reference window of \rrebel{}-std moves
the policy toward the Gibbs tilt of its reward at temperature $\beta\sigma$, so
part (i) describes a converged window (with $\beta\sigma$ in place of $\beta$),
and part (ii) the limit of the mirror-ascent iteration of
Proposition~\ref{prop:mirror}. In either case the policy's expected content
sits at the floor up to a shortfall of slope$/\nu$. With $\nu = 5$, and the
measured curves' envelope slopes of $\wSlopeMin$--$\wSlopeMax$ sentiment points
per content point over the binding floors (Section~\ref{sec:res-control}), that
is a fraction of a content point.

\paragraph{How the curve must be evaluated.}
The point of Eq.~\eqref{eq:curve} averages over prompts drawn from the policy.
Averaging evaluation rows in which one prompt is \emph{sampled} per
$(x,\lmb)$ estimates it without bias. Evaluating the policy's most likely prompt
does so only when the policy is deterministic. Under
Proposition~\ref{prop:eps} the optimal policy is not deterministic, and
evaluating its mode would reinstate the clustering of Lemma~\ref{lem:cluster}.
Our Reward~II evaluation therefore samples prompts from the training sampler
(Section~\ref{sec:setup-exp}).

\paragraph{Why a proportional multiplier.}
Two textbook alternatives to Eq.~\eqref{eq:reward-exp} fail with standardized
rewards. The first is an \emph{exact per-group projection}: choose the content
weight so that the group's own window target meets the floor exactly. Under
standardization a window moves the policy a fixed distance whatever the reward
differences (Proposition~\ref{prop:trust}). With the policy's mass on two
prompts, the projection can only choose which of the two to favour, and it
overshoots. The second is \emph{dual ascent} on a multiplier per context, which
integrates the violation and oscillates against the same fixed-size steps. A
proportional multiplier has neither problem: it is the gradient of the smooth
penalty $F$, so the iteration has the fixed point of
Proposition~\ref{prop:penalty}. It is the proportional term of a PID-Lagrangian
controller, without the integral term. We simulated all three on the real prompt
statistics of the trained Reward~I policies
(Appendix~\ref{app:sim}). With realistic step sizes the proportional rule keeps
$\mathbb{E}[c]$ within \simPfiveCalib{} points of the floor on average (worst
\simPfiveWorst), against \simProjCalib{} for the sampled projection,
\simDualCalib{} for dual ascent (worst \simDualWorst, with sustained
oscillation) and \simExactCalib{} for the exact projection.

\paragraph{What the theory does not cover.}
The results are about idealized, population-level quantities, and four gaps
separate them from the implementation.
(i) The in-group standard deviation uses $G = 16$ samples, and
$\mathbb{E}[1/\hat\sigma] \ne 1/\sigma$. For Gaussian rewards
$\mathbb{E}[\sigma/\hat\sigma] = 1.054$, so standardized targets are about 5\%
larger on average: a slightly lower effective temperature, not a different
target family.
(ii) Within a window the sampler is the moving policy $\pol$, not $\polref$,
so $\sigma_q$ drifts; Propositions~\ref{prop:trust}--\ref{prop:balance} are
exact only at refresh steps.
(iii) Each window runs $150$ Adam steps rather than being solved exactly, and
optimizer state is reset at every scheduler resume (Section~\ref{sec:limitations}),
so Proposition~\ref{prop:mirror} describes an idealized trajectory.
(iv) Proposition~\ref{prop:penalty} describes converged windows of a
nonparametric policy. A $\lmb$-conditioned network shares parameters across
floors, each window runs a finite number of steps, and the multiplier is
estimated from $16$ prompts, so how closely the trained policies meet their
floors is an empirical question. Section~\ref{sec:res-control} measures it.
